\documentclass[aspectratio=169,11pt]{beamer}
\usepackage{../phanes-beamer}
\renewcommand{\ModuleID}{04}
\title{An OpenMC development workflow}
\subtitle{A worked public teaching model, from pin cell to a finite core}
\author{PHANES Initiative}
\institute{The University of Texas at Austin}
\date{September 12, 2026}
\begin{document}
\begin{frame}[plain]
\titlepage
\end{frame}
\begin{frame}[t,fragile]{The engineering task and its declared limits}
\fontsize{11.5}{14}\selectfont
\begin{itemize}\setlength{\itemsep}{2mm}
\item Build an educational UO2/water model in three stages: a reflective pin cell, a reflective 17$\times$17 assembly, and a finite 3$\times$3-assembly core.
\item Compare the finite-core result for 20 cm versus 30 cm of water reflector while holding the fuel model fixed.
\item Deliver geometry plots, validated inputs, convergence evidence, keff with uncertainty, and a relative fission map.
\item This is a deliberately simplified teaching model: no depletion, thermal feedback, boron, control rods, vessel, or licensed design claim.
\end{itemize}
\end{frame}
\begin{frame}[t,fragile]{OpenMC supplies stochastic transport evidence}
\fontsize{11.5}{14}\selectfont
\begin{itemize}\setlength{\itemsep}{2mm}
\item OpenMC tracks sampled neutron histories through a material/geometry model using nuclear data.
\item An eigenvalue run iterates the fission source and estimates multiplication; the output includes sampling uncertainty.
\item Python builds the model and reads results. The transport executable performs the simulation.
\item An agent can create and inspect these artifacts, but cannot replace a missing statepoint with an invented keff.
\end{itemize}
\vfill{\fontsize{6.5}{8}\selectfont\color{UTgray} Sources: omintro. See references and instructor notes.\par}
\end{frame}
\begin{frame}[t,fragile]{Assign the OpenMC work to the development roles}
\fontsize{10.5}{12.5}\selectfont\setlength{\tabcolsep}{4pt}\renewcommand{\arraystretch}{1.15}
\begin{tabular}{>{\raggedright\arraybackslash}p{\dimexpr 0.24\linewidth-3.84pt\relax}>{\raggedright\arraybackslash}p{\dimexpr 0.76\linewidth-12.16pt\relax}}
\toprule
\textbf{Role} & \textbf{Concrete task in this example} \\ \midrule
Planner & Freeze the dimensions, materials, boundaries, outputs, and run budget; identify stages and review gates. \\[2pt]
Researcher & Verify material, lattice, source and statepoint APIs against the selected OpenMC version. \\[2pt]
Implementer & Create model.py, input cases and extraction code without changing the accepted physics spec. \\[2pt]
Verifier / tester & Inspect geometry and XML, test invariants, reproduce a failure, review convergence and statistics. \\[2pt]
Docs / user & Document exact commands and limitations; user accepts geometry and authorizes higher-cost runs. \\[2pt]
\bottomrule\end{tabular}\par\vspace{2mm}
\end{frame}
\begin{frame}[t,fragile]{Start the TUI with an explicit OpenMC request}
\begin{lstlisting}
Implement the supplied public teaching OpenMC specification.
First produce a plan and list unresolved physical assumptions.
Keep nuclear data read-only and do not install packages.
Build pin, assembly and core inputs, plus geometry plots.
Do not start transport until I accept the geometry and budget.
Never report a keff that is not read from the named statepoint.
\end{lstlisting}
\fontsize{10.5}{12.5}\selectfont
\begin{itemize}\setlength{\itemsep}{2mm}
\item The user supplies the approved environment and nuclear-data path.
\item A later approval can authorize a class of bounded smoke runs; each routine command need not be renegotiated.
\end{itemize}
\end{frame}
\begin{frame}[t,fragile]{Freeze the geometric specification}
\fontsize{10.5}{12.5}\selectfont\setlength{\tabcolsep}{4pt}\renewcommand{\arraystretch}{1.15}
\begin{tabular}{>{\raggedright\arraybackslash}p{\dimexpr 0.23\linewidth-3.68pt\relax}>{\raggedright\arraybackslash}p{\dimexpr 0.77\linewidth-12.32pt\relax}}
\toprule
\textbf{Quantity} & \textbf{Teaching value and interpretation} \\ \midrule
Pin pitch & 1.26 cm; square lattice. \\[2pt]
Radii & Fuel 0.4096 cm; gap outer 0.4180 cm; clad outer 0.4750 cm. \\[2pt]
Assembly & 17$\times$17 sites; 25 idealized water sites at index pairs from \{2,5,8,11,14\}. \\[2pt]
Finite core & 3$\times$3 identical assemblies; 200 cm active height. \\[2pt]
Reflector & 20 cm water outside every face, then 30 cm for the comparison. \\[2pt]
\bottomrule\end{tabular}\par\vspace{2mm}
\end{frame}
\begin{frame}[t,fragile]{Make the material model explicit}
\fontsize{10.5}{12.5}\selectfont\setlength{\tabcolsep}{4pt}\renewcommand{\arraystretch}{1.15}
\begin{tabular}{>{\raggedright\arraybackslash}p{\dimexpr 0.2\linewidth-3.2pt\relax}>{\raggedright\arraybackslash}p{\dimexpr 0.8\linewidth-12.8pt\relax}}
\toprule
\textbf{Material} & \textbf{Composition, density and temperature} \\ \midrule
Fuel & UO2; uranium enrichment argument 3 wt\% U-235; O:U atom ratio 2:1; 10.297 g/cm$^3$; 900 K. \\[2pt]
Gap & Helium; 0.0016 g/cm$^3$; 600 K. \\[2pt]
Cladding & Elemental zirconium surrogate; 6.55 g/cm$^3$; 600 K. \\[2pt]
Water & H:O atom ratio 2:1; 0.70 g/cm$^3$; 600 K; H-in-water thermal scattering. \\[2pt]
\bottomrule\end{tabular}\par\vspace{2mm}
\vfill{\fontsize{6.5}{8}\selectfont\color{UTgray} Sources: ommat. See references and instructor notes.\par}
\end{frame}
\begin{frame}[t,fragile]{Construct the fuel and moderator materials}
\begin{lstlisting}
fuel = openmc.Material(name="fuel", temperature=900.0)
fuel.add_element("U", 1.0, enrichment=3.0)
fuel.add_element("O", 2.0)
fuel.set_density("g/cm3", 10.297)
water = openmc.Material(name="water", temperature=600.0)
water.add_element("H", 2.0)
water.add_element("O", 1.0)
water.set_density("g/cm3", 0.70)
water.add_s_alpha_beta("c_H_in_H2O")
\end{lstlisting}
\vfill{\fontsize{6.5}{8}\selectfont\color{UTgray} Sources: ommat. See references and instructor notes.\par}
\end{frame}
\begin{frame}[t,fragile]{Use signed regions for the pin}
\begin{lstlisting}
rf = openmc.ZCylinder(r=0.4096)
rg = openmc.ZCylinder(r=0.4180)
rc = openmc.ZCylinder(r=0.4750)
pin = openmc.Universe(cells=[
    openmc.Cell(fill=fuel, region=-rf),
    openmc.Cell(fill=helium, region=+rf & -rg),
    openmc.Cell(fill=zirconium, region=+rg & -rc),
    openmc.Cell(fill=water, region=+rc)])
\end{lstlisting}
\fontsize{10.5}{12.5}\selectfont
\begin{itemize}\setlength{\itemsep}{2mm}
\item The enclosing root or lattice clips this infinite pin universe.
\item Check that each point belongs to exactly one intended region, away from the ideal mathematical surfaces.
\end{itemize}
\vfill{\fontsize{6.5}{8}\selectfont\color{UTgray} Sources: omgeom. See references and instructor notes.\par}
\end{frame}
\begin{frame}[t,fragile]{Check geometry before running any particles}
\fontsize{11.5}{14}\selectfont
\begin{itemize}\setlength{\itemsep}{2mm}
\item Verify 0 < 0.4096 < 0.4180 < 0.4750 < 1.26/2 cm.
\item The clad-to-clad clearance between neighboring pins is $1.26 - 2(0.4750) = 0.3100$~cm.
\item Assembly width is 17$\times$1.26 = 21.42 cm; finite fuel-region width is 3$\times$21.42 = 64.26 cm.
\item With a 20 cm reflector, outer width is 104.26 cm and outer height is 240 cm.
\end{itemize}
\end{frame}
\begin{frame}[t,fragile]{Build the assembly lattice deliberately}
\begin{lstlisting}
lat = openmc.RectLattice(name="assembly")
lat.pitch = (1.26, 1.26)
lat.lower_left = (-10.71, -10.71)
u = np.full((17, 17), pin, dtype=object)
for row in (2, 5, 8, 11, 14):
    for col in (2, 5, 8, 11, 14):
        u[row, col] = water_universe
lat.universes = u
lat.outer = water_universe
\end{lstlisting}
\vfill{\fontsize{6.5}{8}\selectfont\color{UTgray} Sources: omlattice. See references and instructor notes.\par}
\end{frame}
\begin{frame}[t,fragile]{A finite core changes the boundary problem}
\fontsize{11.5}{14}\selectfont
\begin{itemize}\setlength{\itemsep}{2mm}
\item Fill a 3$\times$3 lattice with the assembly universe and clip it to the 64.26-cm-wide, 200-cm-high fuel-region box.
\item Fill the remainder of the enclosing larger box with water; place vacuum boundaries on the outer box.
\item The pin and assembly stages use reflective boundaries, so their multiplication estimates represent an infinite repetition of the teaching cell.
\item The core stage permits leakage. Do not compare reflective and finite results as if they were the same physical problem.
\end{itemize}
\vfill{\fontsize{6.5}{8}\selectfont\color{UTgray} Sources: omgeom. See references and instructor notes.\par}
\end{frame}
\begin{frame}[t,fragile]{Inspect each stage with a different question}
\fontsize{10.5}{12.5}\selectfont\setlength{\tabcolsep}{4pt}\renewcommand{\arraystretch}{1.15}
\begin{tabular}{>{\raggedright\arraybackslash}p{\dimexpr 0.2\linewidth-3.2pt\relax}>{\raggedright\arraybackslash}p{\dimexpr 0.8\linewidth-12.8pt\relax}}
\toprule
\textbf{Stage} & \textbf{Evidence before promotion} \\ \midrule
Pin & Radial order, fills, pitch, periodic/reflective interpretation and material composition. \\[2pt]
Assembly & 17$\times$17 site count, 25 water sites, intended orientation, outer fill and reflective domain. \\[2pt]
Finite core & Nine assemblies, axial clipping, reflector thickness, vacuum boundaries and x-z plot. \\[2pt]
All stages & No uncovered/overlapping regions in sampled checks; no unexpected material substitutions. \\[2pt]
\bottomrule\end{tabular}\par\vspace{2mm}
\vfill{\fontsize{6.5}{8}\selectfont\color{UTgray} Sources: omtrouble. See references and instructor notes.\par}
\end{frame}
\begin{frame}[t,fragile]{Separate building, plotting, and transport}
\begin{lstlisting}
source .env.openmc            # trusted solver environment
python3 model.py --mode pin --out runs/pin-smoke
cd runs/pin-smoke
openmc --plot
openmc --geometry-debug --threads 4
\end{lstlisting}
\fontsize{10.5}{12.5}\selectfont
\begin{itemize}\setlength{\itemsep}{2mm}
\item model.py writes into a new run directory and refuses to replace an existing one.
\item The instructor provisions the OpenMC executable and OPENMC\_CROSS\_SECTIONS.
\item Transport and geometry-debug execution require the approved compute budget.
\end{itemize}
\vfill{\fontsize{6.5}{8}\selectfont\color{UTgray} Sources: omtrouble. See references and instructor notes.\par}
\end{frame}
\begin{frame}[t,fragile]{Set source and batch controls explicitly}
\begin{lstlisting}
settings.run_mode = "eigenvalue"
settings.particles = 1000
settings.batches = 40
settings.inactive = 10
settings.seed = 3407
settings.source = openmc.IndependentSource(
    space=openmc.stats.Box(lower, upper),
    constraints={"fissionable": True})
\end{lstlisting}
\fontsize{10.5}{12.5}\selectfont
\begin{itemize}\setlength{\itemsep}{2mm}
\item These values are a smoke budget for finding setup defects, not an accepted convergence configuration.
\item The initial source samples the fuel-region box and rejects nonfissionable sites.
\end{itemize}
\vfill{\fontsize{6.5}{8}\selectfont\color{UTgray} Sources: omsource, omsettings. See references and instructor notes.\par}
\end{frame}
\begin{frame}[t,fragile]{Plan convergence runs instead of assuming them}
\fontsize{10.5}{12.5}\selectfont\setlength{\tabcolsep}{4pt}\renewcommand{\arraystretch}{1.15}
\begin{tabular}{>{\raggedright\arraybackslash}p{\dimexpr 0.24\linewidth-3.84pt\relax}>{\raggedright\arraybackslash}p{\dimexpr 0.76\linewidth-12.16pt\relax}}
\toprule
\textbf{Run class} & \textbf{Example budget and purpose} \\ \midrule
Smoke & 1,000 particles $\times$ 40 batches, 10 inactive: inputs/data/geometry and process completion. \\[2pt]
Pilot & 5,000 $\times$ 150 batches, 50 inactive: inspect source behavior and estimate runtime/uncertainty. \\[2pt]
Higher statistics & 20,000 $\times$ 300 batches, 100 inactive: candidate comparison run, subject to convergence review. \\[2pt]
Independent checks & Repeat with different seeds and longer inactive histories if source behavior or comparison warrants it. \\[2pt]
\bottomrule\end{tabular}\par\vspace{2mm}
\vfill{\fontsize{6.5}{8}\selectfont\color{UTgray} Sources: omsettings. See references and instructor notes.\par}
\end{frame}
\begin{frame}[t,fragile]{Source convergence precedes tally acceptance}
\fontsize{11.5}{14}\selectfont
\begin{itemize}\setlength{\itemsep}{2mm}
\item A stable-looking keff trace alone can miss an inadequately converged spatial source.
\item Inspect source diagnostics, including an entropy mesh covering the fission region, and sensitivity to inactive generations.
\item Correlated histories/batches affect uncertainty estimates; investigate suspiciously small reported errors.
\item Only compare reflector cases after both runs meet the predeclared convergence and statistical criteria.
\end{itemize}
\vfill{\fontsize{6.5}{8}\selectfont\color{UTgray} Sources: omsettings. See references and instructor notes.\par}
\end{frame}
\begin{frame}[t,fragile]{Shannon entropy diagnoses fission source convergence}
\fontsize{11.5}{14}\selectfont
\begin{itemize}\setlength{\itemsep}{2mm}
\item Shannon entropy characterizes the spatial distribution of fission sites across batches:
\[H_{\mathrm{src}} = -\sum_{s=1}^{S} P_s \log_2 P_s\]
where $P_s$ is the fraction of fission source sites in spatial mesh cell $s$.
\item The eigenvalue $k_{\mathrm{eff}}$ can appear stable early even while the spatial fission distribution is still settling across the geometry.
\item Inactive batches must run until $H_{\mathrm{src}}$ stationarily fluctuates around its asymptotic mean; only active batches accumulate tallies.
\item Tallying before entropy converges introduces spatial bias and invalidates standard-deviation estimates.
\end{itemize}
\vfill{\fontsize{6.5}{8}\selectfont\color{UTgray} Sources: omsettings. See references and instructor notes.\par}
\end{frame}
\begin{frame}[t,fragile]{Request a spatial fission tally}
\begin{lstlisting}
mesh = openmc.RegularMesh()
mesh.dimension = (51, 51, 1)   # finite teaching core
mesh.lower_left = (-32.13, -32.13, -100.0)
mesh.upper_right = (32.13, 32.13, 100.0)
t = openmc.Tally(name="fission-map")
t.filters = [openmc.MeshFilter(mesh)]
t.scores = ["fission"]
\end{lstlisting}
\fontsize{10.5}{12.5}\selectfont
\begin{itemize}\setlength{\itemsep}{2mm}
\item The mesh resolves pin-pitch columns across the active region and integrates over height.
\item A mesh bin containing a water site should not be interpreted as a fueled pin.
\end{itemize}
\vfill{\fontsize{6.5}{8}\selectfont\color{UTgray} Sources: omtally. See references and instructor notes.\par}
\end{frame}
\begin{frame}[t,fragile]{Normalization is part of the result definition}
\fontsize{11.5}{14}\selectfont
\begin{itemize}\setlength{\itemsep}{2mm}
\item The fission tally estimates reactions per source particle under the eigenvalue normalization.
\item For a relative map, divide a bin's mean by the sum of the mapped fission means and document the mapped region.
\item A power map in watts needs a defensible energy/deposition normalization and declared system power; no such conversion is claimed here.
\item Preserve per-bin standard deviations. Ratio uncertainty needs covariance treatment or a stated approximation.
\end{itemize}
\vfill{\fontsize{6.5}{8}\selectfont\color{UTgray} Sources: omtally. See references and instructor notes.\par}
\end{frame}
\begin{frame}[t,fragile]{Read the exact statepoint into a result record}
\begin{lstlisting}
with openmc.StatePoint(path) as sp:
    k_mean = float(sp.keff.nominal_value)
    k_std = float(sp.keff.std_dev)
    tally = sp.get_tally(name="fission-map")
    means = tally.mean.ravel()
    stds = tally.std_dev.ravel()
\end{lstlisting}
\fontsize{10.5}{12.5}\selectfont
\begin{itemize}\setlength{\itemsep}{2mm}
\item Require the expected statepoint filename in the manifest; do not glob for an arbitrary newest file.
\item Record statepoint hash, input/XML hashes, solver version, exit status and convergence-review status.
\end{itemize}
\vfill{\fontsize{6.5}{8}\selectfont\color{UTgray} Sources: omstate. See references and instructor notes.\par}
\end{frame}
\begin{frame}[t,fragile]{Bind each run to its actual inputs}
\fontsize{10.5}{12.5}\selectfont\setlength{\tabcolsep}{4pt}\renewcommand{\arraystretch}{1.15}
\begin{tabular}{>{\raggedright\arraybackslash}p{\dimexpr 0.23\linewidth-3.68pt\relax}>{\raggedright\arraybackslash}p{\dimexpr 0.77\linewidth-12.32pt\relax}}
\toprule
\textbf{Record} & \textbf{Example field or artifact} \\ \midrule
Input identity & case JSON, model.py hash, exported XML hashes and nuclear-data release identifier. \\[2pt]
Execution & OpenMC version, seed, batch controls, thread/MPI setup, command, return code and log. \\[2pt]
Result & Named statepoint hash, keff mean/std, raw tally arrays and diagnostic plots. \\[2pt]
Acceptance & Geometry review, convergence assessment, precision criterion, reviewer and limitations. \\[2pt]
\bottomrule\end{tabular}\par\vspace{2mm}
\end{frame}
\begin{frame}[t,fragile]{Run the reflector comparison as two cases}
\begin{lstlisting}
python3 model.py --mode core --reflector-cm 20 \
  --particles 20000 --batches 300 --inactive 100 \
  --seed 3407 --out runs/core-r20
python3 model.py --mode core --reflector-cm 30 \
  --particles 20000 --batches 300 --inactive 100 \
  --seed 9181 --out runs/core-r30
\end{lstlisting}
\fontsize{10.5}{12.5}\selectfont
\begin{itemize}\setlength{\itemsep}{2mm}
\item Hold composition, temperatures, geometry inside the fuel region and tally definition fixed.
\item Run only after geometry/budget review; different seeds support the independent-run uncertainty exercise.
\end{itemize}
\end{frame}
\begin{frame}[t,fragile]{Extend the comparison into a bounded sweep}
\fontsize{11.5}{14}\selectfont
\begin{itemize}\setlength{\itemsep}{2mm}
\item Declare the parameter, unit, value set and grid before any run; a sweep is one declared experiment, not repeated single runs.
\item Give each run its own inputs, manifest and result record so the whole grid reproduces from the declaration alone.
\item Fix the statistical target per run first, because comparing under-converged values is not evidence.
\item A partial sweep is reported as partial: name the covered cells, the budget spent, and the remaining grid.
\end{itemize}
\end{frame}
\begin{frame}[t,fragile]{Work the comparison before interpreting it}
\fontsize{11.5}{14}\selectfont
\begin{itemize}\setlength{\itemsep}{2mm}
\item Illustrative values only: case A has $k = 1.00200 \pm 0.00020$; case B has $k = 1.00320 \pm 0.00025$.
\item For independent estimates, $\Delta k = 0.00120$ and $\sigma(\Delta k) = \sqrt{0.00020^2 + 0.00025^2} \approx 0.000320$.
\item Thus $\Delta k \times 10^5 = 120 \pm 32$~pcm in $k$-units; the difference is about $3.75\sigma$.
\item Reactivity $\rho = (k-1)/k$; $\Delta \rho = \Delta k / (k_A k_B) \approx 119 \pm 32$~pcm, showing why reactivity and $\Delta k$ differ.
\end{itemize}
\end{frame}
\begin{frame}[t,fragile]{Choose the next run from the evidence}
\fontsize{11.5}{14}\selectfont
\begin{itemize}\setlength{\itemsep}{2mm}
\item Suppose the target sampling uncertainty for $\Delta k$ is 20 pcm and the current independent-run estimate is 32 pcm.
\item Under the usual inverse-square-root scaling, active histories scale as $(32/20)^2 = 2.56$.
\item Check source convergence and systematic changes before simply increasing histories.
\item Ask the user to authorize the revised budget with the expected cost and precision benefit.
\end{itemize}
\end{frame}
\begin{frame}[t,fragile]{Predeclare checks an agent may not waive}
\fontsize{10.5}{12.5}\selectfont\setlength{\tabcolsep}{4pt}\renewcommand{\arraystretch}{1.15}
\begin{tabular}{>{\raggedright\arraybackslash}p{\dimexpr 0.21\linewidth-3.36pt\relax}>{\raggedright\arraybackslash}p{\dimexpr 0.79\linewidth-12.64pt\relax}}
\toprule
\textbf{Check} & \textbf{Failure and required response} \\ \midrule
Geometry & Overlap, missing region or lost particle: retain diagnostics; fix/review geometry before comparison. \\[2pt]
Data & Missing nuclide or thermal table: report exact requirement; do not silently swap library or material. \\[2pt]
Run identity & Statepoint/input mismatch: reject reuse and rerun the affected case. \\[2pt]
Convergence & Unresolved source behavior: withhold the physical comparison and propose a study. \\[2pt]
Precision & Uncertainty above target: label inconclusive or request the justified additional budget. \\[2pt]
\bottomrule\end{tabular}\par\vspace{2mm}
\end{frame}
\begin{frame}[t,fragile]{Worked failure: the agent removes thermal scattering}
\fontsize{11.5}{14}\selectfont
\begin{itemize}\setlength{\itemsep}{2mm}
\item Observed failure: the configured library cannot supply the requested H-in-water table at the needed conditions.
\item Bad repair: delete the thermal-scattering assignment so transport starts.
\item Verifier finding: this changes the moderator model and invalidates the agreed comparison.
\item Correct next step: identify the missing data/version requirement and ask the operator to supply an approved compatible dataset, or explicitly revise the model and rerun both cases.
\end{itemize}
\vfill{\fontsize{6.5}{8}\selectfont\color{UTgray} Sources: ommat. See references and instructor notes.\par}
\end{frame}
\begin{frame}[t,fragile]{What the user reviews at each OpenMC gate}
\fontsize{11.5}{14}\selectfont
\begin{itemize}\setlength{\itemsep}{2mm}
\item Model gate: dimensions, composition, boundary conditions, simplifications, plots and geometry checks.
\item Run gate: data identity, source controls, planned histories, resource limits and stop conditions.
\item Comparison gate: consistent inputs, accepted convergence, uncertainty, relative tally interpretation and retained diagnostics.
\item Handoff gate: reproducible commands, observed versus illustrative results, unresolved limits and a complete provenance package.
\end{itemize}
\end{frame}
\begin{frame}[t,fragile]{Lab: complete the OpenMC evidence package}
\fontsize{11.5}{14}\selectfont
\begin{itemize}\setlength{\itemsep}{2mm}
\item Use the supplied generator and spec tests to build pin, assembly and finite-core cases in new directories.
\item Inspect plots and geometry diagnostics; document one deliberately introduced geometry or data failure.
\item Execute the approved reflector comparison and replace the illustrative statistics with values extracted from real statepoints.
\item Submit XML, manifests, logs, statepoint references, convergence review, comparison table and a report stating all simplifications.
\end{itemize}
\end{frame}
\begin{frame}[t,fragile]{Sample submission: the OpenMC evidence package}
\fontsize{11.5}{14}\selectfont
\begin{itemize}\setlength{\itemsep}{2mm}
\item Directory layout: pin, assembly, and finite-core cases, each with generated XML, run manifest, logs, and extracted statistics; every statistic shown is taken from your own statepoints, and the illustrative slide values are placeholders.
\item Documented failure: one deliberately introduced geometry or data failure (for example, a dimension that makes two surfaces coincide) with its diagnostic message; the corrected case is rerun and the two manifests are diffed.
\item Comparison table: the reflector and bare-core rows are filled from the extracted k and its standard deviation; the convergence review cites the batch statistics and the relative tally interpretation.
\item Report: all simplifications are stated (public teaching model, no feedback, no depletion), and the handoff commands reproduce every artifact.
\end{itemize}
\end{frame}
\begin{frame}[t,fragile]{Extend the model only through reviewed changes}
\fontsize{11.5}{14}\selectfont
\begin{itemize}\setlength{\itemsep}{2mm}
\item A real assembly may need guide-tube walls, axial regions, burnable absorbers and a documented lattice layout.
\item A reactor study may need heterogeneous loading, depletion, thermal-hydraulic feedback and additional structures.
\item Each extension changes verification, data requirements, state identity and validation evidence.
\item Carry the same planner-to-verifier loop forward; model fidelity and acceptance scope must advance together.
\end{itemize}
\end{frame}
\begin{frame}[t,fragile]{Exit ticket: verifying a transport calculation}
\fontsize{11.5}{14}\selectfont
\begin{itemize}\setlength{\itemsep}{2mm}
\item Why does a stable $k_{\mathrm{eff}}$ trace across early batches fail to prove spatial convergence?
\item How does an engineering agent prove that a reported $k_{\mathrm{eff}}$ originates from transport rather than recall?
\item Why cannot reflective pin-cell and finite-core multiplication values be directly compared?
\item What must happen when a transport run fails because an $S(\alpha,\beta)$ table is unavailable?
\end{itemize}
\end{frame}
\begin{frame}[t,fragile]{A complete answer: transport verification}
\fontsize{11.5}{14}\selectfont
\begin{itemize}\setlength{\itemsep}{2mm}
\item Spatial lag: the fission source distribution converges more slowly than the scalar eigenvalue; Shannon entropy over a fission mesh must plateau before active batches start.
\item Statepoint provenance: verify the run manifest, execution return code, input and XML hashes, and read $k_{\mathrm{eff}}$ directly from the named binary statepoint.
\item Problem domain: reflective pin cells model an infinite lattice without leakage; finite cores include physical leakage and vacuum boundaries.
\item Consequential exception: report the missing library table to the user/operator; do not silently delete thermal scattering or alter moderator physics to force completion.
\end{itemize}
\end{frame}
\begin{frame}[t]{References 1}
\fontsize{9}{12}\selectfont\begin{itemize}\setlength{\itemsep}{3mm}
\item \textbf{curriculum}: User-supplied PHANES curriculum: mod/01 through mod/08/README.md.
\item \textbf{colors}: \href{https://umac.utexas.edu/brand-center/colors/}{umac.utexas.edu/brand-center/colors}
\item \textbf{omintro}: \href{https://docs.openmc.org/en/stable/usersguide/beginners.html}{docs.openmc.org/usersguide/beginners.html}
\item \textbf{ommat}: \href{https://docs.openmc.org/en/stable/usersguide/materials.html}{docs.openmc.org/usersguide/materials.html}
\item \textbf{omgeom}: \href{https://docs.openmc.org/en/stable/usersguide/geometry.html}{docs.openmc.org/usersguide/geometry.html}
\item \textbf{omlattice}: \href{https://docs.openmc.org/en/stable/pythonapi/generated/openmc.RectLattice.html}{docs.openmc.org/pythonapi/generated/openmc.RectLattice.html}
\item \textbf{omtrouble}: \href{https://docs.openmc.org/en/stable/usersguide/troubleshoot.html}{docs.openmc.org/usersguide/troubleshoot.html}
\end{itemize}\end{frame}
\begin{frame}[t]{References 2}
\fontsize{9}{12}\selectfont\begin{itemize}\setlength{\itemsep}{3mm}
\item \textbf{omsource}: \href{https://docs.openmc.org/en/stable/pythonapi/generated/openmc.IndependentSource.html}{docs.openmc.org/pythonapi/generated/openmc.IndependentSource.html}
\item \textbf{omsettings}: \href{https://docs.openmc.org/en/stable/usersguide/settings.html}{docs.openmc.org/usersguide/settings.html}
\item \textbf{omtally}: \href{https://docs.openmc.org/en/stable/usersguide/tallies.html}{docs.openmc.org/usersguide/tallies.html}
\item \textbf{omstate}: \href{https://docs.openmc.org/en/stable/pythonapi/generated/openmc.StatePoint.html}{docs.openmc.org/pythonapi/generated/openmc.StatePoint.html}
\end{itemize}\end{frame}
\end{document}